\documentclass[11pt]{article}
\usepackage[T1]{fontenc}
\usepackage{amsmath,amssymb,amsthm}
\usepackage[margin=1in]{geometry}
\usepackage{xurl}
\usepackage[hidelinks]{hyperref}
\newtheorem{theorem}{Theorem}
\newtheorem{lemma}[theorem]{Lemma}
\newtheorem{proposition}[theorem]{Proposition}
\newcommand{\Spec}{\operatorname{Spec}}
\title{Constructive local--global passage by finite trees\\
\large Certificates, the first-unit scan, and standard \'etale charts}
\author{}
\date{31 October 2026}
\begin{document}
\maketitle

\part{The mathematical argument}

\section{The output determines the useful local statement}

Let $A\to E$ be a map of commutative rings. The desired global
description consists of finitely many $s_i\in E$, coefficients $a_i\in E$,
and standard \'etale presentations of $E[1/s_i]$, with
\[
                 \sum_i s_i a_i=1.
\]
Thus an appropriate local output is a finite algebraic certificate for
one such chart, together with information allowing its denominator to
participate in a finite cover. Merely asserting the existence of a chart
at each existing prime is a different kind of input.

The following construction works intuitionistically. It uses neither
prime-ideal existence nor excluded middle. Its input is an inductively
constructed finite derivation, not an arbitrary pointwise assertion
about primes. Producing that derivation from a particular local proof
is part of the work.

\section{Radical/unit branches}

Fix a ring $R$. A state consists of two finite lists $I$ and $U$ of
elements of $R$. Write $(I)$ for the ideal generated by the first list,
$M(U)$ for the monoid generated by the second, and put
\[
               S(I;U)=M(U)+(I).
\]
Here membership in the sum retains witnesses. This multiplicatively
closed set contains $1$. In $R_{S(I;U)}$ the elements of $U$ are units
and the extension of $(I)$ lies in the Jacobson radical.

A test of an element $a\in R$ produces two successors:
\[
          (I,a;U)\qquad\text{and}\qquad(I;U,a).
\]
The first successor puts $a$ in the radical ideal; the second makes it
a unit. Both successors are retained. This construction does not assert
that one can decide whether $a$ is a unit in the parent ring.

\begin{lemma}[One branching step]\label{lem:branch}
If $x\in S(I,a;U)$ and $y\in S(I;U,a)$, there is
$t\in S(I;U)$ lying in the ideal $(x,y)$.
\end{lemma}
\begin{proof}
Write $x=s+ab$, where $s\in S(I;U)$, and
$y=va^k+j$, where $v\in M(U)$ and $j\in(I)$.
Then
\[
                 t=s^k v+(-b)^k j
\]
belongs to $S(I;U)$. Modulo $(x)$, $s=-ab$, so
$t=(-b)^k y$ modulo $(x)$. Expanding the difference of powers gives
an explicit expression for $t$ in $(x,y)$.
\end{proof}

\begin{theorem}[Certified finite cover]\label{thm:cover}
Start at the empty state and form a finite binary tree of radical/unit
branches. At every leaf $(I;U)$ give an element $s\in S(I;U)$ and a
certificate $C(s)$. Then the leaf elements form a finite family
$s_1,\ldots,s_n$, with all the certificates $C(s_i)$ retained, and there
are coefficients $a_i\in R$ such that $\sum_i s_i a_i=1$.
\end{theorem}
\begin{proof}
Let $J$ be the ideal generated by the finitely many leaf elements.
Every leaf supplies an element of its denominator monoid lying in $J$.
Apply Lemma~\ref{lem:branch} upwards through the tree. At the root,
$S(\varnothing;\varnothing)=\{1\}$, so $1\in J$.
The membership witness supplies the coefficients. Flattening the same
finite tree retains its leaf certificates.
\end{proof}

There is no appeal to an infinite tree or to K\"onig's lemma. Tree
formation, substitution of a finite tree at a leaf, and the proof of
the theorem are finite structural inductions. If witnesses have been
provided only under propositional truncation, the conclusion is also
taken under truncation.

\section{A decision procedure actually used by the local proof}

Consider a monic list of coefficients
\[
                       c_0,c_1,\ldots,c_{d-1},1.
\]
Test $c_0$ by the two successors above. On the unit branch, stop.
On the radical branch, test $c_1$, and continue in this way. After at
most $d$ tests the coefficient $1$ terminates the process. Each leaf
therefore supplies a first unit coefficient, with all earlier
coefficients in the radical ideal of that leaf. This constructs the
tree rather than assuming its existence.

Here is exactly why that information is enough. Let $C\to B$ be a
ring map, let $\mathfrak m\subseteq C$ and $J\subseteq B$ be ideals with
$\mathfrak m B\subseteq J$, and let $s\in B$ satisfy a polynomial
relation modulo $J$. Suppose a finite prefix certificate says that
the earlier coefficients lie in $\mathfrak m$ and the next is a unit.
For a partial Horner expression, the two algebraic steps are
\[
 s^j(c+sH)\in J,\quad c\in\mathfrak m
       \quad\Longrightarrow\quad s^{j+1}H\in J,
\]
and, when $c$ is a unit,
\[
 s^j(c+sH)\in J
       \quad\Longrightarrow\quad s^j(1-sv)\in J,
       \qquad v=-c^{-1}H.
\]
Images of coefficients in $B$ are implicit in these formulas. Reading
the prefix gives $s^k(1-sv)\in J$. Set $N=k+1$ and $e=(sv)^N$.
Then
\[
 e^2-e\in J,\qquad s^N(1-e)\in J,\qquad
 e=s\bigl(s^{N-1}v^N\bigr).
\]
Moreover, if $s^K w\in J$, then $ew\in J$. To see the last assertion,
work modulo $J$: $e$ is idempotent and divisible by $s$, hence
$e=e^K$ is divisible by $s^K$ when $K>0$; the case $K=0$ is immediate.
These are the relation, idempotent and annihilation properties needed
by the local argument. Once the prefix certificate is present, their
proof uses no further decisions.

\section{Actual standard charts}

Use the standard presentation
\[
        E\simeq \bigl(A[X]/(Q)\bigr)_{GQ'}.
\]
This is expressed intrinsically by saying that maps $E\to B$ correspond
to elements $x\in B$ satisfying $Q(x)=0$ and with $(GQ')(x)$ invertible.

\begin{proposition}\label{prop:standard}
If $E$ is standard for $(Q,G)$ and $x$ is its distinguished generator,
then $E[1/F(x)]$ is standard for $(Q,GF)$. The polynomial $Q$ is
unchanged, so monicity is retained when it is given.
\end{proposition}
\begin{proof}
A map from this localisation to $B$ is a map $E\to B$ for which $F(x)$
becomes a unit. The standard universal property identifies this with
an element satisfying $Q(x)=0$ and both $(GQ')(x)$ and $F(x)$ invertible.
Their product is $((GF)Q')(x)$. In a commutative ring a product is a
unit precisely when both factors are units. These correspondences give
the required universal property.
\end{proof}

For a concrete two-branch example, branch on $a=F(x)\in E$.
On the radical branch choose the denominator $1-a=1+(-a)$; on the unit
branch choose $a$. Proposition~\ref{prop:standard} constructs the two
charts, for $(Q,G(1-F))$ and $(Q,GF)$. Theorem~\ref{thm:cover} supplies
the certified finite cover. Neither $a$ nor $1-a$ needs to be a unit
in $E$.

\section{What the constructive bridge still requires}

For the full local structure theorem one must construct a tree whose
leaves contain chart certificates over the \emph{original base}.
The first-unit scan above is one real piece of this construction.
It does not replay the whole field argument, factor selection, ZMT
application, or localisation descent automatically.

In particular, a chart over a node ring must descend, with its finitely
many equations and inverse-map certificates, to a suitable denominator
in that node's $S(I;U)$. Descent from a localised base to a monic
presentation over the original base is also a separate step. Assuming
that a chart exists somewhere in the node localisation is not a
substitute for giving these finite certificates. This is the precise
distinction between the checked finite-tree bridge and the still
unfinished global structure theorem from its original hypotheses.

\part{Connection with the Unit development}

\section{Checked statements}

All new code is in \path{unit/sag/etale-dynamical-global/}. It imports
the existing constructive local and Horrocks libraries by relative
links. No classical segment is imported.

\begin{itemize}
\item \path{LargeTreeCover.seg}: finite trees with leaf data in the
next universe, including maps, substitution, and continuation from the
existing small tree procedures. This permits actual algebras and
universal properties at leaves.
\item \path{DynamicalCertifiedCover.seg}: Theorem~\ref{thm:cover}, with
an explicit finite family, all certificates, and the equation
$\sum_i s_i a_i=1$, under truncation. Only the small denominator data
is sent to the existing \path{FormalMonoidCover}; the larger
certificates are retained by structural recursion. No resizing
principle is needed.
\item \path{DynamicalEtaleCover.seg}: exact chart and localisation
interfaces matching the local theorem, together with a stronger monic
interface. Its main theorem still takes a finite tree of such leaves.
\item \path{FirstUnitData.seg}, \path{FirstUnitRelation.seg}, and
\path{DynamicalCoefficientScan.seg}: the constructed monic scan and
its interpretation in any algebra over the leaf ring, with an ideal
containing the leaf radical ideal. The decision-free algebraic proofs
are extracted from \path{etale-local/IdempotentModM.seg}; that shared
file is unchanged.
\item \path{DynamicalScanCover.seg}: compose that actual scan with
a continuation constructing certified leaves. The scan is proved;
the continuation is an explicit parameter.
\item \path{StandardChartLocalisation.seg}:
Proposition~\ref{prop:standard} with the full algebra universal
property. A presentation and a representative for $F(x)$ construct the
localisation using the existing presentation library.
\item \path{DynamicalComplementaryCover.seg} and
\path{DynamicalStandardCover.seg}: the two constructed branches and
their actual monic standard charts. Polynomial representatives are
chosen only inside the final truncation.
\end{itemize}

\section{Validation and remaining integration}

The fresh development checks 5,789 declarations, including 367 new
declarations in eleven proof modules, in about two seconds. The
\texttt{.opens} inventory covers all 367 new named declarations.
Regression instances include large certificates, the two leaf positions,
constant, linear and quadratic monic scans, a skipped zero coefficient,
the degree-zero case over the zero ring, and actual constructed standard
algebras. Deliberately changing the cover equation, the monic leading
coefficient, the branch direction, or the standard denominator is
rejected by the checker.

The existing univalence, replacement and surjection parameters are
unchanged. There are no new axioms, no checker changes, and no edits
to shared formal sources. Neither a general transformation of arbitrary
local proofs nor the full global \'etale theorem is asserted here.
The next mathematical obligation is to organise the remaining local
steps into compatible finite branch procedures and descend the final
chart certificates; the tree-to-cover part is now available.

\end{document}
